\documentclass{article}
\usepackage[T1]{fontenc}
\usepackage{lmodern}
\usepackage{graphicx}
\usepackage{amsmath,amssymb}
\usepackage[a4paper,margin=2cm]{geometry}
\usepackage[absolute,overlay]{textpos}
\setlength{\TPHorizModule}{1mm}
\setlength{\TPVertModule}{1mm}
\usepackage[indonesian]{babel}
\usepackage{hyperref}
\setlength{\emergencystretch}{2em}
% unit: Halaman 38

\begin{document}

% === Halaman 38 (python/native) ===

38

6. Kunci







Prinsip Kherkoff: semua algoritma kriptografi harus publik (tidak rahasia), 

hanya kunci yang harus rahasia.

•

 Agar enkrpsi dan dekripsi hanya dapat dilakukan oleh dua pihak yang 

berkomunikasi, maka diperlukan kunci rahasia.

• Kunci adalah parameter yang digunakan di dalam enkripsi dan dekripsi 

     Simbol: K    ( K dapat berupa integer, string, alphanumeric, dsb )


Enkripsi: EK(P) = C 

; 

Dekripsi: DK(C) = P

Enkripsi

Plainteks

Dekripsi

Cipherteks

Plainteks

Kunci

Kunci

\end{document}
